\documentclass[../main.tex]{subfiles}

\begin{document}

\section{Chapter 1: Introduction to Data Visualization}
\label{sec:chapter1-introduction}

\subsection{Techniques / Principles Applied}

\begin{itemize}[leftmargin=*]
\item High-dimensional data reduction through visual summaries rather than
row-by-row table inspection.
\item Exploratory and explanatory visualization roles within the same
dashboard.
\item Policy-facing presentation that highlights actionable disparities instead
of isolated statistics.
\end{itemize}

\subsection{How Applied in the Dashboard}

The raw dataset in \texttt{Findex\_Microdata\_2025\_updateViet Nam.csv}
contains 1,000 rows and 199 columns of individual survey information. In raw
tabular form, the links between demographics, digital access, payment
behavior, and resilience are difficult to see quickly. Visualization addresses
that problem by turning survey weights into readable marks such as bar length,
color intensity, and flow width, so the dashboard can show where exclusion and
cash dependence are concentrated.

The dashboard also combines explanatory and exploratory roles. Its fixed
five-page sequence guides the reader from population context to outcomes, while
interactive filters on pages such as Account Ownership \& Access Landscape and
Financial Health \& Resilience let analysts inspect specific subgroups. This
combination is appropriate for a policy audience because it keeps the national
story coherent while still supporting diagnosis.

\subsection{Notes / Adjustments}

Many visualization guidelines separate exploratory tools from guided
presentations, but this dashboard places filters inside an authored narrative.
That deviation is justified because financial inclusion questions are strongly
intersectional, so readers often need to move from the national story to a
specific subgroup without leaving the dashboard.

\end{document}
